% ===================== 封面与使用说明 =====================
\hypersetup{pageanchor=false}
\begin{titlepage}
\centering
\thispagestyle{empty}
\vspace*{2.0cm}

{\color{mainblue}\Huge\bfseries 大学英语六级（CET-6）\par}
\vspace{0.45cm}
{\color{mainblue}\Huge\bfseries 复习指南\par}

\vspace{0.9cm}
{\Large\color{black!70} 核心词汇 · 写译话题词汇与短语 · 真题分类汇编\par}

\vspace{1.6cm}
{\large 目标考试：2026 年 12 月\par}
\vspace{0.3cm}
{\large 适用：SJTU 在校本科生自主备考\par}

\vfill

\begin{minipage}{0.86\textwidth}
\begin{notebox}[title=资料说明]
\begin{itemize}
  \item 本指南为个人备考资料汇总，真题题目、话题与词汇线索整理自公开网络资料（新东方四六级网、新东方在线、英语真题在线等），出处见\hyperref[sec:references]{附录“参考资料”}。
  \item 真题内容仅用于个人学习、练习，版权归原命题机构所有；引用时已尽量精简并标注年份与出处。
  \item 最新考情请以全国大学英语四、六级考试委员会及其官方网站（cet.neea.edu.cn）发布为准。
\end{itemize}
\end{notebox}
\end{minipage}

\vspace{1.0cm}
\end{titlepage}
\hypersetup{pageanchor=true}

\section*{使用说明}
\phantomsection
\addcontentsline{toc}{section}{使用说明}

本指南共分九章，按“先词汇、后真题、再策略”的顺序编排，适合 9--12 月的冲刺复习：

\begin{itemize}
  \item \textbf{第 1 章 考试概况}：题型、分值、时间分配与考场流程，先建立全局感。
  \item \textbf{第 2 章 核心词汇与释义}：按词性整理的高频词表、易混词辨析与固定搭配，适合每天定量默背。
  \item \textbf{第 3 章 写译话题分类词汇与短语}：写作、翻译两大部分的高频话题词汇、搭配与句式，全部按话题分类。
  \item \textbf{第 4--6 章 真题分类汇编}：依次为写作、翻译、选词填空·阅读·听力三个板块的历年真题整理。\textbf{题目均为英文原文呈现}（汉译英部分保留中文原文，因其本身即为翻译任务）。
  \item \textbf{第 7 章 分题型解题策略}：写作模板、翻译技巧、阅读与听力方法、时间分配。
  \item \textbf{第 8 章 备考计划}：以 2026 年 12 月考试为目标的十周计划表。
  \item \textbf{第 9 章 参考资料}：所有素材来源链接。
\end{itemize}

\begin{tipbox}[title=怎么用这份指南]
\begin{itemize}
  \item \textbf{词汇}：第 2 章每天滚动复习 1--2 组；第 3 章的话题词汇在练写译时随时查阅，先“认脸”再“动用”。
  \item \textbf{真题}：第 4--6 章按板块刷题。写作题目自己先计时写，再看第 7 章模板比对；翻译先自己译，再看参考译文。
  \item \textbf{策略}：第 7 章不必通读，哪块错得多先看哪块。
  \item 建议把第 4--6 章打印出来，或对照本 PDF 在纸上作答；真题的“手感”是刷 PDF 之外的网站练不出来的。
\end{itemize}
\end{tipbox}
